\documentclass[11pt]{article}
\usepackage{amsmath,amssymb,amsthm,bm,mathtools}
\usepackage{geometry}
\usepackage{hyperref}
\geometry{margin=1in}

\title{On the Link Between Bayesian Physics-Informed Neural Networks (B-PINNs) and Gaussian Processes (GPs)}
\author{}
\date{}

\newtheorem{proposition}{Proposition}
\newtheorem{corollary}{Corollary}
\newtheorem{remark}{Remark}

\begin{document}
\maketitle

\section*{Scope and Goal}

These notes explain how B-PINNs relate to Gaussian processes in the \emph{infinite-width} regime.
The key message is:

\begin{center}
\emph{In the infinite-width limit, a Bayesian neural network (BNN) prior induces a GP prior on the solution field $u(\cdot)$.\\
If the PDE and boundary operators enter linearly in $u$, then the full B-PINN posterior is exactly the posterior of a GP regression problem with operator observations.}
\end{center}

For nonlinear PDE constraints, the GP prior still holds in the limit, but the posterior is no longer Gaussian; one can then obtain principled Gaussian approximations (e.g.\ Laplace / linearization).

\section{B-PINN in One Line}

Let $u$ solve a PDE on a domain $D$ with boundary $\Gamma$:
\begin{equation}
N_x(u;\lambda)=f \quad (x\in D),\qquad B_x(u;\lambda)=b \quad (x\in \Gamma),
\end{equation}
where $N_x$ is a differential operator and $B_x$ a boundary operator.

A B-PINN represents $u$ by a neural surrogate $\tilde u(x;\theta)$ and defines model-implied
\begin{equation}
\tilde f(x;\theta)=N_x(\tilde u(x;\theta);\lambda),\qquad \tilde b(x;\theta)=B_x(\tilde u(x;\theta);\lambda),
\end{equation}
typically computed via automatic differentiation.

With Gaussian measurement noise, the likelihood is (schematically)
\begin{equation}
P(\mathcal D\mid \theta)\propto \exp\!\Big(
-\tfrac12\| \tilde u(X_u;\theta)-\bar u\|_{\Sigma_u^{-1}}^2
-\tfrac12\| \tilde f(X_f;\theta)-\bar f\|_{\Sigma_f^{-1}}^2
-\tfrac12\| \tilde b(X_b;\theta)-\bar b\|_{\Sigma_b^{-1}}^2
\Big),
\end{equation}
where $X_u,X_f,X_b$ are sensor locations and $\Sigma_u,\Sigma_f,\Sigma_b$ encode noise variances.

The posterior is
\begin{equation}
P(\theta\mid \mathcal D)\propto P(\mathcal D\mid \theta)\,P(\theta),
\qquad
U(\theta):=-\log P(\mathcal D\mid \theta)-\log P(\theta).
\end{equation}

\paragraph{Important conceptual point.}
Even though $U(\theta)$ looks like a PINN loss (sum of squared residuals + prior), B-PINNs aim at
\emph{posterior expectations} such as $\mathbb E[u(x)\mid \mathcal D]$, not a single minimizer (MAP).

\section{Infinite-Width Limit: BNN Prior $\Rightarrow$ GP Prior}

Consider a fully-connected network with width $n$ in hidden layers and i.i.d.\ Gaussian weights/biases
with appropriate scaling (``NNGP'' scaling). Denote the random function
\begin{equation}
u_n(x):=\tilde u(x;\theta),\qquad \theta\sim P(\theta) \text{ (i.i.d.\ Gaussian on weights/biases).}
\end{equation}

\begin{proposition}[NNGP limit (informal statement)]
Fix any finite set of inputs $x_1,\dots,x_m\in \mathbb R^d$.
Under standard assumptions (i.i.d.\ Gaussian initialization, width $n\to\infty$, proper variance scaling),
\begin{equation}
\big(u_n(x_1),\dots,u_n(x_m)\big)\ \xRightarrow[n\to\infty]{}\ \mathcal N(0, K),
\qquad K_{ij}=k(x_i,x_j),
\end{equation}
for some kernel $k$ that depends on the activation function, depth, and prior variances.
Equivalently, $u_n(\cdot)\Rightarrow u_\infty(\cdot)$ where
\begin{equation}
u_\infty(\cdot)\sim \mathcal{GP}(0,k(\cdot,\cdot)).
\end{equation}
\end{proposition}

\begin{remark}[Kernel recursion (conceptual)]
For an MLP, the kernel $k$ can be expressed through a layer-wise recursion:
starting from $k^{(0)}(x,x')$ built from the input inner product, define
$k^{(\ell+1)}(x,x')$ via expectations of $\phi(a)\phi(a')$ for jointly Gaussian $(a,a')$ with covariance
determined by $k^{(\ell)}$.
Closed forms exist for some activations; otherwise the recursion can be evaluated numerically.
\end{remark}

\section{Operator Observations of a GP}

A crucial bridge to physics is: B-PINNs evaluate differential operators on $\tilde u$.
In the GP limit, differential operators applied to a GP remain Gaussian \emph{if the operator is linear}.

Let $u\sim \mathcal{GP}(m,k)$ on $D$ and let $L$ be a linear differential operator (e.g.\ $L=\partial_x$,
$L=\Delta$, or combinations). Under standard smoothness assumptions on $k$ (enough derivatives),
$Lu$ is well-defined as a (mean-square) GP.

\begin{proposition}[Linear operator applied to a GP]
Let $u\sim \mathcal{GP}(m,k)$ and let $L$ be a linear differential operator.
Then $Lu$ is a GP with mean and covariance
\begin{equation}
m_L(x)= (L m)(x),\qquad
k_L(x,x') = (L_x)(L_{x'})\,k(x,x').
\end{equation}
Moreover, the joint process $(u,Lu)$ is Gaussian and the cross-covariance is
\begin{equation}
\mathrm{Cov}\big(u(x),Lu(x')\big) = (L_{x'})\,k(x,x').
\end{equation}
\end{proposition}

\begin{remark}[Boundary operators]
If $B$ is linear (e.g.\ Dirichlet evaluation $u|_\Gamma$, Neumann $n\cdot \nabla u|_\Gamma$),
the same result holds with $L$ replaced by $B$ provided the kernel has the required regularity
up to the boundary.
\end{remark}

\section{Linear PDE Case: GP Regression with Physics Observations}

Assume the physics is \emph{linear} in $u$:
\begin{equation}
Lu = f \quad \text{in } D,\qquad Bu=b \quad \text{on } \Gamma,
\end{equation}
where $L$ and $B$ are linear operators.

Suppose we observe noisy measurements of:
\begin{itemize}
\item $u$ at interior points $X_u=\{x_u^{(i)}\}_{i=1}^{N_u}$,
\item $f=Lu$ at points $X_f=\{x_f^{(i)}\}_{i=1}^{N_f}$,
\item $b=Bu$ at boundary points $X_b=\{x_b^{(i)}\}_{i=1}^{N_b}$,
\end{itemize}
with independent Gaussian noise:
\begin{equation}
\bar u = u(X_u) + \varepsilon_u,\ \varepsilon_u\sim \mathcal N(0,\Sigma_u),\quad
\bar f = (Lu)(X_f) + \varepsilon_f,\ \varepsilon_f\sim \mathcal N(0,\Sigma_f),\quad
\bar b = (Bu)(X_b) + \varepsilon_b,\ \varepsilon_b\sim \mathcal N(0,\Sigma_b).
\end{equation}

Now assume the infinite-width limit prior:
\begin{equation}
u \sim \mathcal{GP}(0,k).
\end{equation}

\subsection{Stacked linear observation model}

Define the stacked observation vector
\begin{equation}
y := \begin{bmatrix}
\bar u\\ \bar f\\ \bar b
\end{bmatrix},
\qquad
\varepsilon := \begin{bmatrix}
\varepsilon_u\\ \varepsilon_f\\ \varepsilon_b
\end{bmatrix},
\qquad
\Sigma := \mathrm{diag}(\Sigma_u,\Sigma_f,\Sigma_b).
\end{equation}

Define the linear observation map (acting on functions $u$)
\begin{equation}
\mathcal A u :=
\begin{bmatrix}
u(X_u)\\ (Lu)(X_f)\\ (Bu)(X_b)
\end{bmatrix}.
\end{equation}

Then the data model is linear-Gaussian:
\begin{equation}
y = \mathcal A u + \varepsilon,\qquad \varepsilon\sim \mathcal N(0,\Sigma).
\end{equation}

\subsection{Closed-form GP posterior}

Because $(u,\mathcal A u)$ is jointly Gaussian, the posterior $u\mid y$ is again a GP:
\begin{equation}
u\mid y \sim \mathcal{GP}\big(m_{\text{post}},\,k_{\text{post}}\big).
\end{equation}

Let $K_{yy}$ denote the covariance of $\mathcal A u$:
\begin{equation}
K_{yy} := \mathrm{Cov}(\mathcal A u,\mathcal A u).
\end{equation}
Using operator covariance identities, $K_{yy}$ can be built blockwise:
\begin{equation}
K_{yy}=
\begin{bmatrix}
K_{uu} & K_{u,Lu} & K_{u,Bu}\\
K_{Lu,u} & K_{Lu,Lu} & K_{Lu,Bu}\\
K_{Bu,u} & K_{Bu,Lu} & K_{Bu,Bu}
\end{bmatrix}
\end{equation}
with entries (example)
\begin{align}
[K_{uu}]_{ij} &= k(x_u^{(i)},x_u^{(j)}),\\
[K_{Lu,Lu}]_{ij} &= (L_x L_{x'})k(x,x')\big|_{x=x_f^{(i)},\,x'=x_f^{(j)}},\\
[K_{u,Lu}]_{ij} &= (L_{x'})k(x,x')\big|_{x=x_u^{(i)},\,x'=x_f^{(j)}},
\end{align}
and similarly for boundary blocks involving $B$.

For any test point $x$, define cross-covariance vector
\begin{equation}
k_{x,y}:=\mathrm{Cov}\big(u(x),\mathcal A u\big)
=
\begin{bmatrix}
k(x,X_u) & (L_{x'})k(x,x')|_{x'=X_f} & (B_{x'})k(x,x')|_{x'=X_b}
\end{bmatrix}.
\end{equation}

Then the posterior mean and covariance are
\begin{equation}
m_{\text{post}}(x) = k_{x,y}\,(K_{yy}+\Sigma)^{-1}\,y,
\end{equation}
\begin{equation}
k_{\text{post}}(x,x') = k(x,x') - k_{x,y}\,(K_{yy}+\Sigma)^{-1}\,k_{y,x'}.
\end{equation}

\begin{corollary}[Equivalence: linear B-PINN (infinite width) $\Leftrightarrow$ GP regression]
Consider a B-PINN with a BNN prior on $u$ and a Gaussian likelihood built from noisy observations
of $u$, $Lu$ (physics), and $Bu$ (boundary). In the infinite-width limit where the BNN prior converges to
$u\sim \mathcal{GP}(0,k)$, if $L$ and $B$ are linear operators then the B-PINN posterior over $u$ is exactly
the GP posterior given by the closed-form formulas above.
\end{corollary}

\paragraph{Interpretation.}
For linear PDEs, the ``physics-informed'' term corresponds to adding \emph{operator observations} (e.g.\ Laplacian,
gradient, boundary evaluation) in a GP regression model.

\section{What Changes for Nonlinear PDEs?}

Suppose physics is nonlinear in $u$, e.g.
\begin{equation}
Lu + g(u) = f,
\end{equation}
with $g$ nonlinear (e.g.\ $g(u)=u(u^2-1)$ as in Allen--Cahn).

In the infinite-width limit:
\begin{itemize}
\item The \emph{prior} over $u$ is still GP: $u\sim \mathcal{GP}(0,k)$.
\item But the likelihood depends nonlinearly on $u$ through $g(u)$.
\item Hence the posterior $P(u\mid \mathcal D)$ is generally \emph{not Gaussian}.
\end{itemize}

\subsection{Principled Gaussian approximations (useful in a project)}

Two standard approximations turn the nonlinear posterior into an approximate GP posterior:

\paragraph{(i) Linearization around a reference function.}
Let $u_\star$ be a reference (e.g.\ MAP / mean). Then
\begin{equation}
g(u)\approx g(u_\star) + g'(u_\star)(u-u_\star),
\end{equation}
turning the nonlinear observation into an \emph{approximately linear} operator observation in $u$.

\paragraph{(ii) Laplace approximation in function space.}
Define the negative log-posterior functional
\begin{equation}
\Phi(u) = -\log P(\mathcal D\mid u) - \log P(u).
\end{equation}
Let $u_{\text{MAP}}$ minimize $\Phi(u)$. Then approximate
\begin{equation}
P(u\mid \mathcal D) \approx \mathcal{GP}\Big(u_{\text{MAP}},\; (\nabla^2\Phi(u_{\text{MAP}}))^{-1}\Big),
\end{equation}
where $\nabla^2\Phi$ denotes the (formal) Hessian operator. In discretized form, this becomes a standard Gaussian approximation.

\section{A Practical Recipe for a Master-Level Theoretical Project}

A rigorous, feasible plan is to focus on \textbf{linear PDEs} and demonstrate the finite-to-infinite width connection.

\subsection*{Theory deliverables}
\begin{enumerate}
\item State the NNGP limit (BNN prior $\Rightarrow$ GP prior).
\item Prove/justify operator covariance identities: $k_{Lu,Lu}(x,x')=L_xL_{x'}k(x,x')$ and cross-covariances.
\item Derive the closed-form GP posterior for mixed observations $(u,Lu,Bu)$.
\item Formulate a clear equivalence statement: infinite-width linear B-PINN $\Leftrightarrow$ GP regression with operator observations.
\end{enumerate}

\subsection*{Experimental deliverables}
\begin{enumerate}
\item Choose a linear PDE benchmark (e.g.\ 1D Poisson $u''=f$ with Dirichlet BC).
\item Empirically estimate the NNGP kernel by sampling many random networks and computing $\hat k(x,x')$.
\item Build the GP posterior using operator kernels (e.g.\ $k_{u'',u''}=\partial_x^2\partial_{x'}^2 k$).
\item Compare to finite-width B-PINN-HMC outputs for increasing width and show convergence toward the GP posterior.
\end{enumerate}

\section{Summary}

\begin{itemize}
\item B-PINNs place a BNN prior on the solution $u(\cdot)$ and incorporate physics through a likelihood on PDE residuals.
\item As width $\to\infty$, the BNN prior converges to a GP prior (NNGP).
\item For linear PDE and boundary operators, ``physics observations'' are linear operator observations of a GP.
\item Therefore, the infinite-width B-PINN posterior is exactly the GP regression posterior with operator observations.
\item For nonlinear PDEs, the prior may still be GP in the infinite-width limit, but the posterior is non-Gaussian; one can use principled Gaussian approximations (linearization, Laplace).
\end{itemize}

\end{document}